% ============================================================================
%  Beyond the Listed Firm
%  A working paper on data infrastructure for AI-entrepreneurship research.
%  Compiles with pdflatex on Overleaf (standard packages only).
% ============================================================================
\documentclass[11pt]{article}

\usepackage[margin=1.1in]{geometry}
\usepackage{amsmath,amssymb,amsthm}
\usepackage{booktabs}
\usepackage{graphicx}
\usepackage{array}
\usepackage{longtable}
\usepackage{caption}
\usepackage[hidelinks]{hyperref}
\usepackage{natbib}
\usepackage{setspace}
\usepackage{enumitem}
\usepackage{footmisc}

\bibliographystyle{plainnat}

\newtheorem{definition}{Definition}
\newtheorem{remark}{Remark}

\title{\textbf{Beyond the Listed Firm:}\\[0.3em]
\large A Data Infrastructure for Measuring the Global AI Startup Ecosystem,\\
and What It Reveals About the Limits of Commercial Startup Data}

\author{Alastair Page \and Jihwan (Victor) Park\thanks{%
Tobin Center for Economic Policy, Yale University. We thank Professor Song Ma
for his guidance and for framing the broader research agenda on artificial
intelligence and the global startup ecosystem. All errors are our own.
Correspondence: \texttt{victorjhp11@gmail.com}.}}

\date{August 2026}

\begin{document}
\maketitle

\begin{abstract}
\noindent
Empirical research on entrepreneurship runs almost entirely on a small number
of commercial company databases. If those databases miss firms
non-randomly---and if what they miss correlates with the technology under
study---then the measured relationship between artificial intelligence and
firm formation is contaminated at the source. We build a pipeline that
collects, extracts, normalizes, resolves, and validates company records from
two commercial databases, public funder registries in the United States and
the European Union, an open-source code host, and 410 heterogeneous
international web sources, and we use it to measure the part of the ecosystem
the commercial databases do not cover. Three results follow. First, the
uncovered layer is not a random remainder: $20.4\%$ of the $52{,}276$
companies we observe that appear in neither commercial database are
AI-related, against $12.3\%$ and $10.6\%$ in the two databases themselves.
The firms that are hardest to see are disproportionately the firms the
research agenda is about. Second, the cost of seeing them is not uniform---
public funder records yield unlisted firms at a $100\%$ rate and cost
nothing, whereas scraping the portfolio pages of well-known venture firms
yields $23$--$30\%$ and costs the most per record. Third, and most
consequential for practice, we document a failure mode of language-model
enrichment: when permitted to estimate fields their input text did not state,
the model filled $7{,}553$ values with its own prior---San Francisco, founded
2020, 11--50 employees---producing a dataset whose descriptive statistics
would have been publishable and wrong. Because every value was written with a
source and a confidence, the imputed values were identified and withdrawn
exactly. We argue that per-value provenance is not bookkeeping but the
enabling condition for using language models in dataset construction at all.
We close by demonstrating what the infrastructure supports: a panel of
$641{,}442$ firms observed in 2021 and 2025, used to date and decompose the
adoption of AI identity, and a set of research questions the resulting data
can and cannot answer.
\end{abstract}

\medskip
\noindent\textbf{Keywords:} entrepreneurship measurement, artificial
intelligence, data infrastructure, coverage bias, entity resolution,
LLM-assisted extraction

\vspace{1em}
\noindent\textbf{Note on data.} This paper reports aggregates only. No
row-level records from any licensed commercial dataset are reproduced or
redistributed here. The two commercial startup databases used are referred to
as \emph{Database~A} and \emph{Database~B} throughout.

\newpage
\tableofcontents
\newpage

% ============================================================================
\section{Introduction}
% ============================================================================

The question motivating the broader research program this paper belongs to is
whether artificial intelligence is changing how firms are formed, financed,
and staffed around the world---and whether it lowers the barriers to founding
a company or entrenches the advantages of incumbents, elite founders, and
established clusters. That question is empirical. Answering it requires
observing the population of new firms.

The practical difficulty is that nobody observes that population. What
researchers observe is a \emph{register}: a commercial database assembled by a
vendor whose incentives are commercial rather than statistical. A firm enters
such a register when someone has a reason to enter it---when it raises
institutional capital, when it is acquired, when a journalist writes about it,
when a founder submits a profile. A firm that has a product, a website, a
codebase, customers, and federal grant funding, but no venture round and no
press coverage, may simply never appear.

This is a familiar measurement problem, but AI gives it a specific and
uncomfortable shape. The firms most likely to be missing are young, small,
technical, software-first, and not yet venture-financed. Those are also, on
priors, the firms most likely to be AI-native. If that is true, then the
measured share of AI among new firms is biased downward by an amount that is
itself a function of the technology---and every downstream estimate of AI's
effect on entrepreneurship inherits the bias.

\subsection{What this paper does}

We do not attempt to estimate the causal effect of AI on firm formation. We
build the instrument that such an estimate would require, characterize what
it can and cannot see, and use it to measure the coverage gap directly.

Concretely, the paper makes four contributions.

\paragraph{1. A formalized, operating pipeline.}
Section~\ref{sec:pipeline} specifies the system as a composition of stages---
collection, extraction, normalization, entity resolution, validation---each
with an explicit policy and an explicit failure mode. The system is not a
prototype: it comprises roughly $43{,}000$ lines of Python across a collection
backend, an analysis layer, and a public dashboard, developed over $189$
commits between February and August 2026, and it holds on the order of
$9.9\times10^{5}$ company records. What distinguishes it from a large scraper
is that collection strategy is chosen \emph{per source} by a documented
policy, and every stage writes an audit record.

\paragraph{2. A measurement of the observability frontier.}
Section~\ref{sec:frontier} defines the \emph{unlisted layer}---firms our
pipeline observes that appear in neither commercial database---and shows that
it differs systematically from the covered population along AI intensity,
sector, vintage, and geography. The AI-adoption gap ($20.4\%$ vs.\ $12.3\%$
and $10.6\%$) is the paper's central empirical claim, and it is the one claim
that requires no enrichment and no imputation: it is computable for
$100\%$ of the population from data present at collection time.

\paragraph{3. A negative result about LLM enrichment, and the discipline that
caught it.}
Section~\ref{sec:enrichment} reports what happened when we allowed a language
model to estimate company attributes its input text did not state. The model
did not decline. It returned the modal startup---San Francisco, founded 2020,
11--50 employees---for company after company, at rates ($33.7\%$, $19.1\%$,
$70.6\%$) that would have plotted as findings. We report the episode in full
because the recovery is the methodological point: since every field was
written with a source tag and a confidence, all $7{,}553$ estimated values
were located and removed without disturbing anything else. We state this as a
design rule for LLM-assisted dataset construction: \emph{the model may read,
but it may not recall}.

\paragraph{4. A demonstration of what the infrastructure enables.}
Section~\ref{sec:longitudinal} uses the assembled panels to date the adoption
of AI identity among established firms, decompose genuine repositioning from
relabeling, and link both to firm outcomes. We present these as descriptive
demonstrations with their identification problems stated, not as causal
estimates. Section~\ref{sec:limits} is a full accounting of what the data
cannot support.

\subsection{Relation to the broader research agenda}

This paper is one technical and empirical component of a larger program
studying how artificial intelligence is reshaping global entrepreneurship,
venture financing, and innovation. Its role in that program is to establish
what can be measured, at what cost, with what error, and with what residual
blind spots---so that later work on formation, financing, founder
composition, and geographic clustering rests on a characterized instrument
rather than an assumed one.

% ============================================================================
\section{The Measurement Problem, Stated Formally}
\label{sec:problem}
% ============================================================================

Let $\mathcal{F}$ denote the true population of firms founded in some
window---an object no one observes. A data source $s$ is a subset
$\mathcal{O}_s \subseteq \mathcal{F}$ together with a map from firms to
records. Research is conducted on $\mathcal{O}_s$ but reasons about
$\mathcal{F}$.

Define the \emph{inclusion propensity} of firm $i$ in source $s$:
\begin{equation}
\pi_s(i) \;=\; \Pr\!\left[\, i \in \mathcal{O}_s \mid \mathbf{z}_i \,\right],
\end{equation}
where $\mathbf{z}_i$ collects firm characteristics. The standard implicit
assumption in empirical entrepreneurship is that $\pi_s(i)$ is either constant
or independent of the outcome of interest. Write $y_i^{AI}\in\{0,1\}$ for
whether firm $i$ is AI-related. The naive estimator of AI prevalence from
source $s$ is
\begin{equation}
\widehat{p}_s \;=\; \frac{1}{|\mathcal{O}_s|}\sum_{i \in \mathcal{O}_s} y_i^{AI},
\end{equation}
and its bias relative to the population share $p = \mathbb{E}[y^{AI}]$ is
governed by the covariance between inclusion and the outcome:
\begin{equation}
\label{eq:bias}
\mathbb{E}[\widehat{p}_s] - p \;\;\propto\;\; \mathrm{Cov}\!\left(\pi_s(i),\, y_i^{AI}\right).
\end{equation}

Equation~\eqref{eq:bias} is the object of this paper. It cannot be signed from
theory, and it cannot be estimated from $\mathcal{O}_s$ alone---nothing inside
a database reveals what the database omits. It can, however, be bounded from
below by constructing a second observation channel $\mathcal{O}_{\text{web}}$
with a different inclusion mechanism and comparing $y^{AI}$ across
$\mathcal{O}_{\text{web}} \setminus \mathcal{O}_s$ and $\mathcal{O}_s$. That
comparison is what Section~\ref{sec:frontier} reports, and building
$\mathcal{O}_{\text{web}}$ is what Sections~\ref{sec:pipeline}--\ref{sec:enrichment}
describe.

\begin{remark}[What a coverage comparison can and cannot establish]
Our web channel is a convenience sample, not a probability sample: there is no
sampling frame and therefore no denominator. Consequently we make no claim
about the level of $p$. The claim is about a \emph{difference} between
populations we observe, which is identified without a denominator. This
distinction recurs throughout the paper and disciplines which results we are
willing to headline.
\end{remark}

% ============================================================================
\section{The Pipeline}
\label{sec:pipeline}
% ============================================================================

\subsection{Overview and notation}

Let $S$ index sources. Source $s$ emits raw observations
$R_s = \{r_{s1},\dots,r_{sn_s}\}$, where a raw observation may be an HTML
listing row, a JSON object from a backend API, a row in a bulk file, or a code
repository with its metadata. An extraction operator $E_s$ maps raw
observations to a common record schema:
\begin{equation}
E_s(r_{sj}) \;\longmapsto\; \mathbf{x}_{sj} =
\big(\text{name},\, \text{domain},\, \text{description},\,
\text{country},\, \text{city},\, \text{founded},\, \text{batch},\, \dots\big).
\end{equation}
The pooled raw corpus is $\mathcal{R} = \bigcup_{s=1}^{S} R_s$, and the system
implements
\begin{equation}
\mathcal{R}
\;\xrightarrow{\;E\;}\; \mathcal{X}
\;\xrightarrow{\;N\;}\; \tilde{\mathcal{X}}
\;\xrightarrow{\;\rho\;}\; \mathcal{E}
\;\xrightarrow{\;V\;}\; \mathcal{D},
\end{equation}
where $N$ normalizes, $\rho$ resolves records to entities, $V$ validates, and
$\mathcal{D}$ is the analysis dataset. Each arrow is a place where records are
lost or corrupted, and each is instrumented.

\subsection{The source lattice}

Sources differ along two axes that turn out to matter more than any other
design consideration: \emph{marginal cost per record} and \emph{marginal
novelty}---the probability that a record from that source is not already
held. Table~\ref{tab:sources} summarizes the classes.

\begin{table}[t]
\centering
\small
\caption{Source classes and their role in the pipeline.}
\label{tab:sources}
\begin{tabular}{@{}p{3.4cm}p{4.5cm}rp{3.6cm}@{}}
\toprule
\textbf{Class} & \textbf{Instances} & \textbf{Records} & \textbf{Access} \\
\midrule
Commercial database A & startup/company register & $650{,}200$ & licensed bulk file \\
Commercial database B & venture-financing register & $281{,}395$ & licensed bulk files \\
Public funder registries & U.S. federal small-business awards; national health and science agencies; E.U. framework programme & $\sim$$17{,}000$ firms & open bulk download \\
Open-source code host & topic- and keyword-scoped repository search & $\sim$$8{,}000$ orgs & public API \\
Heterogeneous web & accelerators, university venture programmes, venture portfolios, launch directories, national ecosystem sites & $410$ sites attempted & scraping / agentic \\
Employment database & derived firm- and person-level aggregates & $335{,}842$ rows & licensed \\
\bottomrule
\end{tabular}
\end{table}

The web tier is the most technically demanding and the least productive per
unit of effort; the funder registries are the reverse. This asymmetry is a
finding in its own right and is quantified in
Section~\ref{sec:yield}.

\subsection{Collection policy}

For each source the system selects a collection strategy according to
\eqref{eq:policy}:
\begin{equation}
\label{eq:policy}
C_s \;=\;
\begin{cases}
C^{\text{bulk}} & \text{licensed or open bulk file},\\[2pt]
C^{\text{api}}  & \text{documented backend endpoint (search, REST, index API)},\\[2pt]
C^{\text{rule}} & \text{stable DOM; a deterministic parser exists},\\[2pt]
C^{\text{agent}} & \text{unknown, dynamic, or broken structure},\\[2pt]
C^{\text{excl}} & \text{repeated failure; suspended}.
\end{cases}
\end{equation}

$C^{\text{rule}}$ is implemented by $36$ deterministic scrapers, each a
subclass of a common base that returns validated records and performs no
database writes of its own; persistence, deduplication, and validation are
handled once, in the base class, so that a scraper cannot corrupt the store by
being written carelessly.

$C^{\text{agent}}$ is a tool-using language-model agent with a bounded budget:
at most $10$ reasoning iterations, $12$ plain page fetches, and $8$
browser-rendered fetches per site. It is given three tools---a hosted
extraction service that returns readable text from a URL, a plain HTTP
fetcher, and a headless browser for client-rendered pages---and must decide
which pages to visit, whether pagination exists, and when it has finished.

The policy is not static. Transitions are driven by observed outcomes:

\begin{equation}
\label{eq:transitions}
\begin{aligned}
C^{\text{rule}} &\to C^{\text{agent}} && \text{after two consecutive failures or a zero-result run},\\
C^{\text{agent}} &\to C^{\text{excl}} && \text{after three consecutive failures},\\
C^{\text{excl}} &\to C^{\text{agent}} && \text{after } 90 \text{ days},\\
\text{any} &\to \text{cooldown} && \text{for } 7 \text{ days after a successful run}.
\end{aligned}
\end{equation}

This is what makes the system maintainable at $410$ sites by two people. Web
sources break constantly and silently; the failure of a scraper is not an
exception to be handled but the steady state to be budgeted for.

\subsection{Instruction caching: the agent as a compiler}

When the agent succeeds on a site it writes a YAML instruction file recording
the seed URLs, the strategy that worked, pagination and subpage hints, and the
quality statistics of the successful run. Subsequent runs read the file and
attempt the cached strategy first, falling back to the full agentic path only
if it fails. The agent thus amortizes: the expensive path is paid once per
site, not once per run, and the artifact it produces is human-readable and
editable.

\begin{table}[t]
\centering
\small
\caption{Agentic collection outcomes across all attempted sites
($N=410$ instruction files, as committed).}
\label{tab:agentic}
\begin{tabular}{@{}lr@{}}
\toprule
Sites attempted (instruction file written) & $410$ \\
Sites with at least one successful extraction & $220$ \quad ($53.7\%$) \\
\midrule
\multicolumn{2}{@{}l}{\emph{Successful strategy, conditional on success}}\\
\quad single-page extraction & $193$ \quad ($87.7\%$) \\
\quad subpage discovery & $15$ \quad ($6.8\%$) \\
\quad pagination probe & $12$ \quad ($5.5\%$) \\
\midrule
\multicolumn{2}{@{}l}{\emph{Yield, conditional on success}}\\
\quad total records & $5{,}040$ \\
\quad mean per site & $22.9$ \\
\quad median per site & $6$ \\
\quad 90th percentile / maximum & $50$ \; / \; $814$ \\
\midrule
\multicolumn{2}{@{}l}{\emph{Validation statistics of the accepted runs}}\\
\quad valid-name ratio (mean) & $1.000$ \\
\quad duplicate ratio (mean) & $0.004$ \\
\quad completeness score (mean) & $0.996$ \\
\bottomrule
\end{tabular}
\end{table}

Table~\ref{tab:agentic} reports the outcome. Three features deserve comment.

First, the success rate is $53.7\%$. An instruction file is written on
\emph{attempt}, not on success, so the file count is a record of what was
tried rather than a library of what works---a distinction we initially got
wrong, and which inflated our own internal projections of remaining yield by
roughly an order of magnitude until the files were audited individually. We
report it because it is the kind of error that is easy to make and invisible
from aggregate counts.

Second, the yield distribution is extremely skewed: mean $22.9$, median $6$.
A small number of large directories dominate the total, and the typical
international site contributes a handful of firms. Any plan that projects
total yield from an average is wrong by a factor of roughly four.

Third, conditional on acceptance, quality is high---near-zero duplicate rate
and near-complete records---because the validator gates acceptance. This is
selection, not accuracy: it tells us that what passes is clean, not that the
extractor is correct on what it rejects. We return to this in
Section~\ref{sec:limits}.

\subsection{Normalization}

$N$ applies deterministic transformations before any matching occurs:

\begin{itemize}[leftmargin=1.4em,itemsep=1pt]
\item \textbf{Domain canonicalization.} URLs are lowercased, stripped of
scheme, \texttt{www.}, path, and query, and reduced to the registrable domain
via public-suffix extraction, so that \texttt{https://www.Example.co.uk/about}
and \texttt{example.co.uk} resolve identically.
\item \textbf{Product-domain filtering.} A blocklist of code hosts,
documentation hosts, social platforms, package registries, and preprint
servers prevents \texttt{github.io} or \texttt{medium.com} from being adopted
as a company's identity---a failure that would silently merge thousands of
distinct firms into one entity.
\item \textbf{Name normalization.} Case folding, punctuation stripping, and
removal of legal suffixes.
\item \textbf{Country normalization.} A curated map to a fixed country list;
values that are not countries (U.S. state names, bare city names, stray
coordinates) resolve to \emph{no country} rather than being coerced. They are
$0.8\%$ of rows with a country value and are excluded from regional counts.
\item \textbf{Incumbent exclusion.} A shared denylist removes large
technology incumbents at import time across every source, so that the largest
firms in the world do not appear in a dataset of emerging startups.
\end{itemize}

\begin{remark}[Refusing to default]
The system stores \texttt{NULL} rather than a plausible value wherever a value
is unknown---no default country, no default city, no imputed founding year in
the primary field. This costs coverage and is the correct trade: a defaulted
value is indistinguishable from a measured one downstream, and defaults
concentrate on exactly the modal category, manufacturing spurious
concentration. Section~\ref{sec:prior} shows what happens when this rule is
relaxed.
\end{remark}

\subsection{Entity resolution}
\label{sec:er}

The same firm appears in a licensed bulk file, on an accelerator's portfolio
page, in a federal award record, and as an organization on a code host, under
four different names. Resolution (stage $\rho$, Section~\ref{sec:er}) proceeds by a strict, auditable
rule rather than a learned score.

Define the entity key
\begin{equation}
\kappa(\mathbf{x}) =
\begin{cases}
\texttt{domain:}\,d(\mathbf{x}) & \text{if a canonical registrable domain } d(\mathbf{x}) \text{ exists},\\
\texttt{name:}\,\nu(\mathbf{x}) & \text{otherwise, with } \nu \text{ the normalized name}.
\end{cases}
\end{equation}
Domain equality is treated as identity, enforced by a unique index. When no
domain is available the system falls back to a similarity rule,
\begin{equation}
\label{eq:match}
\mathrm{Match}(i,j) \;=\; \mathbb{1}\!\left[\, S_n(i,j) \ge \tau \;\wedge\; \Sigma(i,j) \,\right],
\qquad
\tau = 0.92 \;\; (\tau = 0.95 \text{ for database B}),
\end{equation}
where $S_n$ is a normalized string similarity and $\Sigma(i,j)$ requires a
corroborating shared signal (a common organization login, a common outbound
link, a common portfolio listing). Every accepted match is written to an audit
table recording the method used, so that any downstream result can be
recomputed under a stricter rule.

\begin{remark}[Why not a learned matcher]
Equation~\eqref{eq:match} is deliberately conservative and deliberately
unlearned. A probabilistic matcher of the form
$M(i,j)=w_nS_n+w_dS_d+w_lS_l+w_gS_g$ with a fitted threshold would raise
recall, and we sketch it in Section~\ref{sec:future} as the natural next
step---but fitting it requires labeled match pairs we do not have, and the
asymmetry of errors favors precision: a false merge destroys two firms and is
nearly undetectable afterwards, while a false split leaves two recoverable
records. We record the current rule honestly as a rule, and do not describe it
as a model.
\end{remark}

\subsection{Validation and audit}

$V$ accepts a run only if the extracted batch passes a joint quality test on
name validity, duplicate rate, and field completeness (the realized values
appear in Table~\ref{tab:agentic}). Independently of acceptance, every scrape
attempt writes a row recording source, tier, status, counts found and new,
duration, and diagnosis. Per-domain health is tracked separately and drives
the transitions in \eqref{eq:transitions}. The audit tables are what make
statements like ``$53.7\%$ of attempted sites ever succeeded'' computable at
all; a pipeline that logs only successes cannot measure its own coverage.

% ============================================================================
\section{Identifying AI Firms}
\label{sec:classification}
% ============================================================================

Every claim in this paper depends on a binary label $y_i^{AI}$. We describe how
it is produced, and---more usefully---the two ways our first attempts at it
failed.

\subsection{The cascade}

Classification runs cheapest-first:
\begin{equation}
y_i^{AI} =
\begin{cases}
1 & \text{strong AI term in name, description, or repository topics},\\
0 & \text{no technology-adjacent term anywhere},\\
f_\theta(d_i, p) & \text{otherwise (ambiguous): a language model reads the text}.
\end{cases}
\end{equation}
The first branch is high precision and free; the second is a cheap negative
that removes the large majority of records from consideration; only the
ambiguous residual reaches a model call. Terms are matched with word-boundary
constraints, and the bare token ``ai'' is handled separately because it is a
substring of ordinary English and of many proper nouns.

A weighted score is also computed for firms with repository metadata:
\begin{equation}
\label{eq:aiscore}
a_i \;=\; \min\Big(1,\;
0.3\,\mathbb{1}[\text{strong topic}] +
0.2\,\mathbb{1}[\text{moderate topic}] +
0.2\,\mathbb{1}[\text{strong text}] +
0.1\,\mathbb{1}[\text{moderate text}] +
0.2\,\mathbb{1}[\text{database-A AI tag}]\Big),
\end{equation}
with the operative threshold $a_i \ge 0.5$.

\subsection{Failure one: a score that does not transfer across sources}

Equation~\eqref{eq:aiscore} was designed for repository data, where topics are
developer-chosen, high-precision labels. Applied to the commercial databases it
breaks in a specific and instructive way. Records from database~B have no
topics and short descriptions; their attainable maximum is $0.2 + 0.1 = 0.3$,
which is \emph{below the threshold by construction}. A scoring rule calibrated
on one source had, in effect, silently excluded an entire other source of
$281{,}395$ firms.

The obvious repair---lower the threshold---fails a validity check. Setting
$a_i \ge 0.1$ admits roughly $22{,}000$ additional firms, but their founding
years are close to uniform at about $1{,}200$ per year from 2010 to 2025.
Genuine AI-era formation is sharply increasing in time; a flat cohort profile
is the signature of legacy analytics and machine-learning-adjacent firms for
which AI is peripheral. We therefore retained the threshold and adopted an
explicit disjunction over source-specific evidence,
\begin{equation}
y_i^{AI} = \mathbb{1}\big[\,a_i \ge 0.5 \;\vee\; \text{tag}^{A}_i \;\vee\; \ell_i \,\big],
\end{equation}
where $\text{tag}^A_i$ is database~A's own AI taxonomy flag and $\ell_i$ is a
model verdict on the question ``is AI the core product of this company?''
obtained for the ambiguous residual.

The general lesson is that a classifier's threshold is a property of the
\emph{source}, not of the concept. Any cross-source AI indicator must either
be recalibrated per source or defined on evidence that all sources actually
carry.

\subsection{Failure two: a category tag is not a description}
\label{sec:tagtext}

Commercial databases assign category tags. It is tempting to treat a vendor's
``artificial intelligence'' tag as ground truth. We tested this by taking
firms whose category tags changed to include AI between two vintages, and
asking a model to compare their \emph{actual descriptions} across the same
two vintages ($n=500$, complete parse rate). Table~\ref{tab:tagtext} reports the
verdicts.

\begin{table}[h]
\centering
\small
\caption{Do firms newly tagged AI show AI in their own text?}
\label{tab:tagtext}
\begin{tabular}{@{}lrr@{}}
\toprule
Verdict from the description pair & Count & Share \\
\midrule
Genuine pivot to AI & $99$ & $19.8\%$ \\
AI added as a feature & $74$ & $14.8\%$ \\
Already AI in the earlier text & $63$ & $12.6\%$ \\
Marketing language only & $4$ & $0.8\%$ \\
\textbf{No AI in the text at all} & $\mathbf{260}$ & $\mathbf{52.0\%}$ \\
\bottomrule
\end{tabular}
\end{table}

Slightly over half of firms newly tagged as AI show no AI language in their own
self-description in either vintage. Taken at face value, the tag overstates
genuine AI adoption by roughly a factor of three ($34.6\%$ genuine versus
$100\%$ implied by the tag). Explicit text-level ``AI-washing''---marketing
claims with no substance behind them---is rare ($0.8\%$); the dominant error is
not firms lying, but vendor tags moving for reasons unrelated to what the firm
says it does.

\begin{remark}
The practical rule we adopt: \emph{always validate a category tag against
description text before using it as a treatment or an outcome}. This is cheap
and, in our data, changes the headline by a factor of three.
\end{remark}

% ============================================================================
\section{Enrichment Under an Evidence Hierarchy}
\label{sec:enrichment}
% ============================================================================

Collection yields sparse records: often a name and a URL. Enrichment fills
fields. It is also where a pipeline built around language models can quietly
destroy itself, so we specify the rule first and report the violation second.

\subsection{The rule}

Order evidence by whether it can be checked:
\begin{equation}
\label{eq:tiers}
\underbrace{T_1: \text{official registries}}_{\$0}
\;\succ\;
\underbrace{T_2: \text{documents we fetch}}_{\sim\$0.003}
\;\succ\;
\underbrace{T_3: \text{model reads a text we hold}}_{\sim\$0.0007}
\;\succ\;
\underbrace{T_4: \text{model recalls a fact}}_{\text{prohibited}}.
\end{equation}

$T_1$ comprises federal and national award records, an E.U. framework
programme registry, domain registration records, and a code host's
organization API. These are free, checkable, and---critically---were
\emph{recorded for a reason unrelated to our study}. $T_2$ is the firm's own
website, fetched and read. $T_3$ is a model performing an inference over text
we possess and can re-read. $T_4$---asking a model what it knows about a
company---is excluded by construction.

The reason $T_4$ is excluded is not a general objection to language models. It
is specific to this population. These firms are, by definition, the ones no
index covers. A model asked to recall facts about them has nothing to recall,
and a model asked a direct question answers it anyway.

$T_1$ alone was decisive. Founding-year coverage on the unlisted population
rose from $8.6\%$ to $59.5\%$ in a single day using domain-registration dates,
first-award years, and code-host account-creation dates, at zero marginal
cost; $7{,}420$ firms gained named executives from award records. These are
facts a model would otherwise have been asked to supply, and would have
supplied.

\subsection{The violation, and what it produced}
\label{sec:prior}

We relaxed the rule once. In an early enrichment pass the model was permitted
to \emph{estimate} location, founding year, team size, and stage when the
source text did not state them, with the estimate flagged as inferred. It
filled nearly everything. Table~\ref{tab:prior} reports the result.

\begin{table}[h]
\centering
\small
\caption{Output of permitting estimation of unstated fields.}
\label{tab:prior}
\begin{tabular}{@{}lllr@{}}
\toprule
Field & Modal value & Share at mode & Flagged inferred \\
\midrule
City & San Francisco & $33.7\%$ & --- \\
Team size & 11--50 & $70.6\%$ & $99.8\%$ \\
Country & United States & $46.1\%$ & $93.8\%$ \\
Founding year & 2020 & $19.1\%$ & $95.9\%$ \\
\bottomrule
\end{tabular}
\end{table}

This is not noise and it is not hallucination in the usual sense of a
fabricated specific. It is the model's prior over ``startup''---the median
startup, repeated. Plotted, it would have looked like a finding: a
concentration of AI activity in San Francisco, a founding surge in 2020, a
modal team size consistent with seed-stage firms. Each of those is a claim
someone in this literature would want to make, and each would have been an
artifact of the imputation step.

Formally, permitting $T_4$ replaces the unknown quantity with a draw from the
model's marginal:
\begin{equation}
\hat{x}_{ik} \;\sim\; q_\theta(x_k) \quad \text{when } x_{ik} \text{ is unstated},
\end{equation}
independent of $i$. Any statistic computed over the filled column then
converges to a functional of $q_\theta$ rather than of the data. Because
$q_\theta$ is concentrated on the modal startup, the resulting distributions
are \emph{more} concentrated and \emph{more} plausible than the truth, which
is precisely why the error is hard to catch by inspection.

\subsection{Why provenance is the enabling condition}

Every enriched value is written to a JSONB provenance map as
$\{\text{field} \mapsto (\text{source}, \text{confidence})\}$. Formally each
stored value is a triple
\begin{equation}
v_{ik} = \big(x_{ik},\; s_{ik},\; c_{ik}\big),
\qquad s_{ik} \in \{T_1, T_2, T_3, T_4\},
\end{equation}
and a query may filter on $s_{ik}$. Because the estimated values carried
$s_{ik} = T_4$, all $7{,}553$ of them were located and deleted exactly,
together with their provenance entries, leaving every $T_1$--$T_3$ value
untouched. Without per-value provenance the only remedies would have been to
discard the entire enrichment pass or to keep a contaminated dataset; both
were avoided at the cost of one JSONB column.

We state the general principle plainly, because it generalizes well beyond
this project: \textbf{if a pipeline writes model output into the same columns
as measured values without recording which is which, the dataset cannot be
repaired---only rebuilt.}

\subsection{What honest coverage looks like}

The result of enforcing \eqref{eq:tiers} is coverage that is high where records
exist and low where they do not, rather than uniformly high and partly
fictional. For field $k$, define
\begin{equation}
\mathrm{Cov}_k = \frac{1}{N}\sum_{i=1}^{N}\mathbb{1}\big[x_{ik} \ne \varnothing\big].
\end{equation}

\begin{table}[h]
\centering
\small
\caption{Field coverage over $45{,}812$ verified unlisted entities.}
\label{tab:coverage}
\begin{tabular}{@{}lr@{\hskip 2.2em}lr@{}}
\toprule
Field & $\mathrm{Cov}_k$ & Field & $\mathrm{Cov}_k$ \\
\midrule
Founding year (incl.\ proxies) & $73.8\%$ & City & $39.0\%$ \\
Description & $65.8\%$ & Founders / executives & $16.2\%$ \\
AI application, model, sector & $58.9\%$ & Product status & $9.2\%$ \\
Country & $52.6\%$ & Team size & $5.8\%$ \\
Domain & $47.5\%$ & & \\
\bottomrule
\end{tabular}
\end{table}

Table~\ref{tab:coverage} is the number usually reported. It is also the number
that misleads, because fields are not independent and a regression needs them
jointly. Define the \emph{intersection coverage} over a required set $K$:
\begin{equation}
\mathrm{Cov}_K = \frac{1}{N}\sum_{i=1}^{N}\prod_{k \in K}\mathbb{1}\big[x_{ik} \ne \varnothing\big].
\end{equation}
Requiring description, domain, founding year, country, and AI classification
\emph{together} leaves $5{,}778$ firms---$21.4\%$ of the $26{,}981$ classified
AI firms, against per-field figures of $47$--$74\%$. Any specification needing
all five runs on that subsample, and we report it up front rather than in a
footnote.

\subsection{Founding year is a proxy, and says so}

For the unlisted population, a stated founding year is usually unavailable.
The system stores an ordered set of proxies in separate columns, never
overwriting the primary field:
\begin{equation}
\widehat{\text{founded}}_i = \mathrm{COALESCE}\big(
\underbrace{\text{founded}_i}_{\text{stated}},\;
\underbrace{\text{cohort}_i}_{\text{accelerator batch}},\;
\underbrace{\text{grant}_i}_{\text{first award}},\;
\underbrace{\text{web}_i}_{\text{first archived capture}}\big).
\end{equation}
Each successive term is a weaker upper bound on the true founding date: a firm
is founded, then joins an accelerator; is founded, then wins a first federal
award; registers a domain at or after founding, and a reused domain can
predate the firm entirely. Decoding accelerator batch labels alone
(``W26'', ``Spring 2026'', ``2019'') recovered $2{,}414$ years; first-award
extraction recovered $7{,}681$. Because the terms are stored separately, any
analysis can restrict to stated years and observe how much the result moves.

% ============================================================================
\section{The Observability Frontier}
\label{sec:frontier}
% ============================================================================

\subsection{Defining the unlisted layer}

\begin{definition}[Unlisted firm]
A firm in our data is \emph{unlisted} if it appears in neither commercial
database. The \emph{strict} definition additionally requires no match to the
employment database, operationalized as the absence of the industry code that
only that match can set.
\end{definition}

Counting is not the first step; defining what counts is. Most unlisted records
arrive from a code-host scan, where a bare name establishes nothing. Rather
than assume, we asked the code host directly; Table~\ref{tab:funnel} reports the
resulting funnel.

\begin{table}[h]
\centering
\small
\caption{From raw records to an analyzable population.}
\label{tab:funnel}
\begin{tabular}{@{}lrp{6.1cm}@{}}
\toprule
Layer & Count & Basis \\
\midrule
Everything found (unlisted) & $56{,}981$ & --- \\
$-$ personal accounts & $-11{,}169$ & code-host API returned \texttt{type: User} \\
\textbf{Verified entities} & $\mathbf{45{,}812}$ & organization, funder record, or portfolio listing \\
\quad of which analyzable & $36{,}534$ & has a domain or a description \\
\quad \textbf{of which AI firms} & $\mathbf{26{,}981}$ & classified from its own description \\
\bottomrule
\end{tabular}
\end{table}

A name containing a space or a period cannot be a code-host login, so $4{,}527$
records were settled by rule at zero API cost---which also preserves the
meaning of a negative lookup for the names that remained. The intermediate
layers are reported as what they are rather than quietly dropped or quietly
counted.

\subsection{The central result: coverage bias is AI-correlated}

\begin{table}[h]
\centering
\small
\caption{AI share by coverage status. Computable for $100\%$ of each
population; requires no enrichment.}
\label{tab:headline}
\begin{tabular}{@{}lrrr@{}}
\toprule
Population & Firms & AI firms & AI share \\
\midrule
Unlisted, but present in the employment database & $2{,}455$ & $645$ & $26.3\%$ \\
\textbf{Unlisted (strict)} & $\mathbf{52{,}276}$ & $\mathbf{10{,}675}$ & $\mathbf{20.4\%}$ \\
Commercial database A & $650{,}200$ & $80{,}277$ & $12.3\%$ \\
Commercial database B & $281{,}395$ & $29{,}953$ & $10.6\%$ \\
\bottomrule
\end{tabular}
\end{table}

Table~\ref{tab:headline} is the paper's principal empirical claim. Firms that
the commercial registers do not cover are AI-related at roughly $1.7$ to
$1.9$ times the rate of firms they do cover. In the notation of
Section~\ref{sec:problem}, $\mathrm{Cov}(\pi_s, y^{AI}) < 0$: inclusion
propensity and AI status are negatively associated, so
$\widehat{p}_s$ understates $p$.

We emphasize what makes this specific comparison robust where other results in
this paper are not. It requires no founding year, no country, no team size, no
imputation, and no proxy. It is a relative comparison between observed
populations, so the absence of a sampling frame---fatal for level
estimates---does not bind. And it survives the composition problem discussed
below, because AI classification is available for the whole population while
the enriched fields are not.

\subsection{Structure of the unlisted layer}

\paragraph{Sector.} The unlisted layer over-represents healthcare
($10.1\%$ vs.\ $3.2\%$ in database~A), financial services ($9.5\%$ vs.\
$3.3\%$), education, data and analytics, and energy; it under-represents
professional services ($4.7\%$ vs.\ $13.8\%$) and biotechnology. The pattern
is consistent with a simple mechanism: firms with a website and a public
codebase but no financing event or filing history are visible to us and
invisible to the registers, while professional-services and biotech firms
generate exactly the paperwork the registers ingest.

\paragraph{Vintage.} The unlisted layer skews sharply to 2019--2025, with
$10$--$12\%$ of the bucket in each recent year, while database~A's cohort
counts fall away after 2021. Part of this is genuine---the web channel sees
firms before a register indexes them---and part is register entry lag, which
we cannot separate.

\paragraph{Geography.} Of $27{,}588$ unlisted firms with a usable country
value, $57.9\%$ are in North America, $13.2\%$ Western Europe, $7.6\%$ East
Asia, $4.2\%$ Southern Europe, $4.1\%$ South Asia, and $12.2\%$ across nine
other regions; $0.8\%$ carry a value that is not a country and are excluded
rather than assigned. By continent: Americas $59.5\%$, Europe $21.6\%$, Asia
$15.1\%$, Africa $1.8\%$, Oceania $1.2\%$.

\begin{remark}[The composition caveat, stated where it bites]
The geographic and vintage figures are dominated by a specific
sub-population. Of unlisted firms with an effective founding year, roughly
$64\%$ are U.S.\ federal-grant recipients---firms that are $100\%$ U.S.\ by
construction and entered our data through a bulk registry download rather than
through web discovery. They pull the unlisted U.S.\ share from about $40\%$ to
about $60\%$. Any geography or formation-timing result on the unlisted
population is really a result about American deep-tech grant recipients. The
AI-share result in Table~\ref{tab:headline} is the exception precisely because
it needs neither year nor country.
\end{remark}

\paragraph{Activity mix.} A bottom-up taxonomy built from description
embeddings rather than an imposed vertical scheme places the unlisted layer
disproportionately in developer tools and infrastructure ($2{,}450$ of
$31{,}377$ firms in that domain), with the covered population dominating
customer-facing categories. This is the same story as the sector tilt, told
from the product side: the unlisted layer is where the technical supply chain
of AI sits.

\subsection{Discovery yield: not all sources cost the same}
\label{sec:yield}

Because the pipeline records provenance per firm, we can measure each
channel's \emph{novelty}: the share of firms from that channel that turn out
to be in neither commercial register.

\begin{table}[h]
\centering
\small
\caption{Novelty and precision by discovery channel.}
\label{tab:yield}
\begin{tabular}{@{}lrr@{}}
\toprule
Channel & Unlisted yield & Share that are real firms \\
\midrule
Public funder records (federal and E.U.) & $\mathbf{100\%}$ & $99.6\%$ \\
Accelerators and university programmes & $96.5\%$ & --- \\
Lesser-known venture portfolios & $76\%$ & $89.5\%$ \\
Well-known venture portfolios & $23$--$30\%$ & --- \\
Code-host scan & --- & $75.8\%$ \\
\bottomrule
\end{tabular}
\end{table}

The ordering in Table~\ref{tab:yield} is almost exactly inverse to the
ordering of engineering effort. A famous venture firm's portfolio page is
JavaScript-rendered, paginated, frequently redesigned, and requires the
agentic tier---and three-quarters of what it returns is already held, because
a famous firm's portfolio is indexed everywhere. Public funder records are
static bulk downloads requiring no scraping at all, and every firm in them is
new to us and almost certainly a real company.

\begin{remark}[A practical prior for researchers building similar data]
Prioritize bulk-downloadable administrative and funder records over web
scraping. The scraping is more interesting and produces almost nothing that
the registers do not already have. We spent a large share of our effort
learning this.
\end{remark}

\subsection{The formation trend, with its caveat attached}

Across the pooled dataset, the AI share of newly founded firms rises from
$3.0\%$ in 2000 to $9.1\%$ in 2014, $17.3\%$ in 2020, $43.4\%$ in 2023, and
$58.3\%$ in 2025. The direction is unambiguous and consistent across
subpopulations. The level in the final years is not trustworthy: register
entry lag means recent cohorts are observed mainly through their most
visible---and disproportionately AI---members, and the denominator for
2024--2025 is small and still filling. We report the series as a description
of what is observable over time, not as an estimate of AI's share of firm
formation.

% ============================================================================
\section{What the Infrastructure Enables}
\label{sec:longitudinal}
% ============================================================================

The preceding sections describe an instrument. This section demonstrates the
class of question it makes answerable, using panels assembled from multiple
vintages of the commercial registers. We label these demonstrations, not
findings: each is descriptive, and each carries an identification problem we
state rather than solve.

\subsection{A firm-level panel}

Joining a 2021 vintage of database~B to its 2025 vintage on a stable internal
identifier---an exact join, no fuzzy matching---yields $641{,}442$ firms
observed at both dates. AI language prevalence in self-descriptions rises from
$3.8\%$ to $4.7\%$; $7{,}648$ firms add AI language, $2{,}206$ drop it.

Applying a model as a \emph{difference} classifier---given the 2021 and 2025
descriptions of the same firm, what changed?---on $n=2{,}000$ firms yields:
genuine pivot $18.1\%$, AI added as a feature $34.5\%$, already AI in 2021
$10.5\%$, marketing language only $3.4\%$, no real change $33.5\%$. The
decomposition is stable from $n=150$ to $n=2{,}000$.

\paragraph{Triangulating the marketing-only category.} Text alone cannot
separate a firm that repositioned from a firm that rewrote its copy. Using
headcount and financing changes over the same window as an independent
signal, $28\%$ of firms that added AI language show no growth of either kind,
against the $3.4\%$ the text classifier flags. Text \emph{under}-detects
unsupported claims by roughly a factor of eight. Absence of growth is a strong
signal, not proof---a genuine pivot may not yet have grown---but the gap
between $3.4\%$ and $28\%$ is too large to attribute to timing.

\paragraph{Heterogeneity.} Information technology repositions around AI most
often ($2.5\%$ of firms add AI language); B2B services shows the highest
unsupported-claim rate ($20\%$ of its additions). By vintage, the rate is
monotone in recency: pre-2010 $0.7\%$, 2010--15 $1.6\%$, 2016--20 $3.0\%$,
2021+ $4.0\%$. Younger firms reposition roughly four times as often as
pre-2010 firms.

\subsection{Dating AI identity, and a correction we had to make}

Using four consecutive vintages of database~A, $21.3\%$ of firms holding an AI
category tag in the latest vintage already held one in the earliest, $6.0\%$
acquired it in between, and $57.5\%$ are new to the register in the latest
vintage.

The $6.0\%$ figure carries the tag-versus-text problem of
Section~\ref{sec:tagtext} on top of everything below.
The $57.5\%$ is largely an artifact: the latest export is substantially
larger and uses an expanded taxonomy, so ``new to the register'' conflates
formation with vendor coverage growth. More importantly, the $6.0\%$
repackaging figure is a \emph{floor} produced by an inadequate baseline. The
earliest available vintage of database~A postdates the public release of
consumer generative AI, so firms that adopted an AI identity in the
intervening period appear as ``always AI''. Two checks establish this: within
the ``already AI'' pool only $0.3\%$ used generative-AI vocabulary in their
earlier description---they were traditional machine-learning firms---and
bridging to the 2021 vintage of database~B shows that roughly $50\%$ of
today's AI-tagged firms were not AI in 2021.

We report the correction rather than the original because the mechanism is
general and easy to reproduce: \textbf{a panel's earliest vintage silently
becomes the definition of ``always''}, and if that vintage postdates the
treatment, the treatment becomes invisible.

A second correction concerns vocabulary. An early count of firms adopting
``agent'' language was restricted three ways at once---to firms with
descriptions in both panel vintages, to newly added terms, and by a narrow
pattern---returning $598$ firms. Counted across all $137{,}857$ AI firms with
a description in the latest vintage, $13{,}906$ ($10.1\%$) mention agents and
$8{,}257$ ($6.0\%$) use explicitly agentic language---more than mention
generative AI ($3.9\%$), large language models ($2.5\%$), or specific model
families ($1.3\%$). The direction of the original finding held; its magnitude
was understated by more than an order of magnitude. Both corrections are
recorded in the repository alongside the original numbers.

\begin{remark}[On a substring bug]
An earlier version of this analysis reported generative-AI vocabulary at
$0.48\%$ in 2021. The pattern matched the substring \texttt{llm}, which occurs
inside ordinary words such as \emph{enrollment}. Corrected, clean
generative-AI phrasing moves $0.007\% \to 0.011\% \to 0.079\%$ across the three
vintages---essentially zero before the technology existed, which is the
expected result and a useful check that the corrected measure is measuring
what it claims.
\end{remark}

\subsection{Outcomes}

Matching to a register vintage carrying operating status yields $577{,}921$
matched firms. Within founding cohort---which removes the age confound that
makes older non-AI firms appear more successful---firms that added AI language
close at substantially lower rates than same-vintage firms that never did:
$6.5\%$ vs.\ $13.7\%$ (pre-2010), $3.5\%$ vs.\ $17.8\%$ (2010--15), $2.2\%$
vs.\ $6.6\%$ (2016--21). Acquisition and public-listing rates are comparable
within cohort.

This association must not be read causally, for a reason internal to the
measurement: \emph{to add AI language by 2025, a firm had to be alive enough
in 2025 to update its profile}. Survival is partly built into the treatment
definition, which mechanically depresses the closure rate of the treated
group. A credible design needs the date AI was added and a baseline
conditioned on being alive at the start. We report the association and the
obstruction together.

\subsection{Who founds AI firms, and one unreconciled contradiction}

From register-derived founder records, founders of AI firms are more
credentialed and more technical than founders of non-AI firms: $20.5\%$ hold
doctorates versus $10.8\%$; $33.4\%$ hold only a bachelor's degree versus
$47.1\%$; teams are larger ($1.72$ vs.\ $1.39$ founders). Fields concentrate in
computer science, electrical engineering, physics, and mathematics. AI
founders are markedly \emph{less} likely to be serial founders ($3.1\%$ vs.\
$11.9\%$), consistent with a wave of first-time entrepreneurship.

On financing progression, AI firms entering at seed reach the next round at
$23.2\%$ versus $18.8\%$, and the gap persists in every founding cohort.

One measure does not reconcile. The register-derived founder data show AI
founders slightly less likely to be female ($10.0\%$ vs.\ $13.4\%$), while an
employment-database-derived aggregate over a different population shows no
gender difference. The two use different sources, different populations, and
different definitions of ``founder''---the employment source has no founder
role category at all, so its aggregate describes senior officers near
founding, not verified founders. We report the discrepancy and decline to cite
either figure until it is resolved.

\subsection{An informative negative}

For the unlisted population no register panel exists, so we attempted to
reconstruct homepage text at two dates from a web archive. Roughly $88$--$90\%$
of sampled unlisted AI firms have no archived 2021 homepage: they did not
exist, were not online, or were not crawled. The measurement failed, and the
failure is the result---the unlisted AI layer is predominantly born-AI rather
than repositioned, which is consistent with its 2019--2025 vintage skew and its
elevated AI share. Repositioning is a phenomenon of established, covered
firms. We report this because a null that explains itself is worth more than a
tool that quietly returns a biased $12\%$ subsample.

% ============================================================================
\section{Limitations}
\label{sec:limits}
% ============================================================================

We state the limitations at length because the value of this paper is in the
characterization of the instrument, and an instrument characterized only by
its successes is not characterized.

\paragraph{No sampling frame.} The web-discovered population is whatever the
scrapers and scans found. There is no denominator, so no rate is ``per
capita'' and no level is a population estimate. All defensible claims are
relative comparisons between observed populations.

\paragraph{Missingness on the key variable is non-random, and correlated with
the outcome.} Only $8.3\%$ of unlisted firms carry a stated founding year.
Worse, the subsample that does is systematically different on exactly what we
measure: unlisted firms with a founding year are $27.5\%$ AI versus $20.1\%$
without; with a domain, $28.0\%$ versus $16.8\%$. Enrichable firms are more
AI-intensive than unenrichable ones, so any statistic computed on the enriched
subsample \emph{overstates} AI adoption. This is selection on observables and
it biases the enrichment-dependent results in a known direction.

\paragraph{Coverage is uneven across channels, confounding channel
comparisons.} Founding-year coverage is $20.1\%$ for scraper-sourced firms and
$4.1\%$ for code-host-sourced firms, while the code host contributes the
larger population. Channel comparisons on enriched fields compare different
things.

\paragraph{Proxies are bounds, not measurements.} Domain registration and
first award year bound founding from above and are not incorporation dates.
Named people from award records are executives at award time; a principal
investigator is not necessarily a founder, and the job title is stored with
the name for that reason.

\paragraph{Survival is a proxy.} Domain liveness records whether a domain
answered an HTTP request on a given day. Parked domains read as live, lapsed
but operating firms read as dead, and the roughly two-thirds of the population
with no domain cannot be checked. It is directional and is never reported as a
survival rate.

\paragraph{No economic depth for the unlisted layer.} There is no financing,
headcount, or verified exit data for these firms. ``Economic analysis'' of the
unlisted layer means formation timing, sector mix, a survival proxy, and an AI
flag---structure, not performance.

\paragraph{Composition.} The federal-grant sub-population is $100\%$ U.S.\ and
dominates the enriched unlisted subsample; healthcare is likely overstated at
$20.8\%$ of classified firms because a national health agency is one of our
largest registry sources. Activity mix should be read by channel.

\paragraph{Extraction accuracy is not measured against a labeled gold
standard.} We report validation statistics conditional on acceptance and a
small hand-built regression suite for the difference classifier, but we do not
have manually labeled ground truth for field-level extraction. Precision,
recall, and $F_1$ per field are therefore \emph{not} reported, and we decline
to estimate them from the pipeline's own outputs. Building that labeled set is
the highest-value next step and is described below.

\paragraph{Classifier granularity.} The technique sub-classification returns
``other'' for $63.7\%$ of firms; the descriptions are frequently too thin to
identify a method, and the taxonomy likely needs revision.

\paragraph{Operational fragility.} The system depends on third-party APIs with
rate limits, credit exhaustion, and unstable connections; one month's run log
shows a $73\%$ error rate driven overwhelmingly by upstream extraction
failures and rate limiting rather than by parsing errors. Pipelines of this
kind are not run-and-forget, and effort estimates should assume they are not.

% ============================================================================
\section{Next Steps}
\label{sec:future}
% ============================================================================

\paragraph{A labeled evaluation set.} Hand-label a stratified sample of a few
hundred firms across sources and fields, and report per-field precision,
recall, and $F_1$ for the extraction operator $E$, along with a comparison of
classification strategies---keyword, embedding similarity, model, and
hybrid---on the same labels:
\begin{equation}
\text{Accuracy} = \frac{1}{N}\sum_{i=1}^N \mathbb{1}\big[\hat{y}_i = y_i\big],
\qquad
F_1^{(k)} = \frac{2 \,\mathrm{Prec}_k \mathrm{Rec}_k}{\mathrm{Prec}_k + \mathrm{Rec}_k}.
\end{equation}
This is the single largest gap between what the system does and what the paper
can claim about it.

\paragraph{A learned entity resolver.} Replace \eqref{eq:match} with
$M(i,j) = w_n S_n + w_d S_d + w_l S_l + w_g S_g$ over name, domain, location,
and description similarity, with weights fitted on the labeled pairs above and
$\tau$ chosen on a precision-constrained frontier. The audit table already
stores the method used for every existing match, so the current rule can serve
as a weak-supervision seed and any change can be evaluated against it.

\paragraph{A source-reliability model.} With multiple sources per firm, source
quality can be learned from observed agreement rather than assumed:
$R_s = \alpha C_s + \beta A_s + \gamma U_s$ over completeness, cross-source
agreement, and update reliability. We flag this explicitly as \emph{not
implemented}.

\paragraph{Event-study timing.} The repositioning results need the date AI was
added, not two endpoints. A mid-window register vintage plus the archive-based
month-dating already prototyped would support an event study with a
still-alive-at-baseline sample, which is the minimum for a credible design.

\paragraph{Coverage against an external reference.} Where a national business
register provides counts by country and year, the coverage ratio
\begin{equation}
\mathrm{CoverageRatio}_c = \frac{P_{\text{ours}}(c)}{P_{\text{reference}}(c)}
\end{equation}
would convert our relative statements into calibrated ones for those
countries, and would let us reweight toward a population quantity rather than
only comparing observed groups.

% ============================================================================
\section{Conclusion}
% ============================================================================

Empirical work on AI and entrepreneurship inherits the coverage properties of
whatever database it runs on. We built a second observation channel and used
it to measure those properties directly. The firms the commercial registers
miss are AI-related at roughly $1.7$--$1.9$ times the rate of the firms they
cover; the cheapest way to see those firms is administrative funder records
rather than the web scraping that consumed most of our engineering effort; and
the most dangerous step in a modern data pipeline is the one where a language
model is asked to supply a value rather than to read one, because it will
comply, plausibly, at scale, and in a direction that looks like a result.

None of this settles whether AI lowers the barriers to founding a firm or
raises them. It establishes that the question cannot be answered from a single
commercial register without a correction whose sign we can now state and whose
rough magnitude we can now bound, and it provides an instrument, a provenance
discipline, and an honest account of the instrument's blind spots for the work
that follows.

\vspace{1em}
\noindent\textbf{Reproducibility.} The pipeline, analysis scripts, and
aggregate outputs are in the project repository. Licensed row-level data are
excluded by configuration; only aggregates are committed. Every figure in this
paper is regenerable from a committed script against the production database,
and each enriched value carries its source and confidence in the store.

\end{document}
